\chapter{The Quantum Threat}
\label{ch:quantum}

\section{Why this chapter matters}

Every other chapter in this book builds toward schemes labelled \emph{post-quantum}. This chapter explains what that word is reacting to. Post-quantum cryptography exists because a sufficiently large quantum computer would break almost all of the public-key cryptography deployed today -- the very algorithms protecting web traffic, software updates, and stored data right now. Understanding \emph{which} problems fall, \emph{which} survive, and \emph{why} is the motivation for everything that follows.

\section{Learning goals}

After reading this chapter, you should be able to:

\begin{itemize}
  \item name the hard problems that today's public-key cryptography relies on;
  \item explain, at a high level, what Shor's and Grover's algorithms do;
  \item distinguish what a quantum computer \emph{breaks} from what it merely \emph{weakens};
  \item explain ``harvest now, decrypt later'' and why migration is urgent;
  \item place the published NIST standards and the main mathematical families in context.
\end{itemize}

\section{What today's cryptography rests on}

Classical public-key cryptography stands on two number-theoretic problems that are believed hard for ordinary computers:

\begin{itemize}
  \item \textbf{Integer factorization} -- given a large $N = p\cdot q$, recover the primes $p$ and $q$. This is the foundation of \textbf{RSA}.
  \item \textbf{Discrete logarithm} -- given $g$ and $g^x$ in a suitable group, recover $x$. This underlies \textbf{Diffie--Hellman} key exchange and elliptic-curve schemes such as \textbf{ECDSA} and \textbf{ECDH}.
\end{itemize}

For a classical computer, the best known algorithms for both problems run in super-polynomial (sub-exponential or worse) time, so the parameters can be chosen to make an attack take longer than the age of the universe. That safety margin is exactly what a quantum computer removes.

\section{Shor's algorithm: the break}

In 1994 Peter Shor showed that a quantum computer can factor integers and compute discrete logarithms in \emph{polynomial} time~\cite{shor1997}. The engine is \emph{period finding}: both problems can be recast as finding the period of a cleverly chosen function, and a quantum computer extracts that period efficiently using the quantum Fourier transform -- evaluating the function on a superposition of all inputs at once, then interfering the results so that the period stands out.

The consequence is stark: RSA, Diffie--Hellman, ECDSA, and ECDH are \emph{all} broken by a large enough quantum computer, no matter how big the keys are. Doubling an RSA key only adds polynomial work for Shor's algorithm -- it does not restore security. There is no parameter choice that saves them.

\section{Grover's algorithm: the weakening}

Symmetric cryptography (block ciphers like AES) and hash functions (like SHA-3) face a milder threat. Grover's algorithm~\cite{grover1996} speeds up \emph{generic search}: finding a marked item among $N$ possibilities takes about $\sqrt{N}$ steps instead of $N$. Against a symmetric key of $b$ bits, this cuts the effective security to roughly $b/2$ bits.

Crucially, this is a \emph{quadratic} speedup, not an exponential one. For exhaustive key search, AES-256 therefore retains roughly $128$ bits of generic quantum-search security. Hash properties must be treated separately: for an $m$-bit ideal hash, preimage and second-preimage search have a generic Grover-style cost of about $2^{m/2}$, whereas collision search has different quantum algorithms and cost models. In the quantum-query model, the Brassard--H\o yer--Tapp algorithm finds collisions in about $2^{m/3}$ queries; memory, circuit cost, and concrete security estimates require separate analysis. Symmetric primitives \emph{survive}, but their parameters and security goals must be chosen carefully.

\begin{figure}[H]
\centering
\begin{tikzpicture}[node distance=4mm and 14mm, every node/.style={font=\scriptsize}]
  % broken column
  \node[keynode] (rsa)  {RSA};
  \node[keynode, below=of rsa] (dh)  {Diffie--Hellman};
  \node[keynode, below=of dh]  (ec)  {ECDSA / ECDH};
  \node[draw=pqcorange, thick, rounded corners=2pt, fill=pqclight,
        right=22mm of dh, align=center, font=\small] (shor) {Shor\\[-1pt]\scriptsize (polynomial time)};
  \node[right=20mm of shor, pqcorange, font=\bfseries] (broken) {BROKEN};
  \draw[flowarrow, pqcorange] (rsa.east) -- (shor);
  \draw[flowarrow, pqcorange] (dh.east) -- (shor);
  \draw[flowarrow, pqcorange] (ec.east) -- (shor);
  \draw[flowarrow, pqcorange] (shor) -- (broken);
  % survive column
  \node[flownode, below=16mm of ec] (aes)  {AES, SHA-3};
  \node[flownode, below=of aes] (pqc) {lattice / hash / code problems};
  \node[draw=pqcteal, thick, rounded corners=2pt, fill=pqclight,
        right=22mm of aes, yshift=-6mm, align=center, font=\small] (grov) {Grover\\[-1pt]\scriptsize (only $\sqrt{N}$)};
  \node[right=20mm of grov, pqcteal, font=\bfseries] (survive) {SURVIVE};
  \draw[flowarrow, pqcteal] (aes.east) -- (grov);
  \draw[flowarrow, pqcteal] (pqc.east) -- (grov);
  \draw[flowarrow, pqcteal] (grov) -- (survive);
\end{tikzpicture}
\caption{The quantum threat, sorted. Shor's algorithm breaks the number-theoretic public-key schemes outright (orange). Grover's algorithm gives generic quadratic speedups for search (teal); its implications for keys, preimages, and collisions are not identical.}
\label{fig:quantum-threat}
\end{figure}

\begin{table}[H]
\centering
\begin{tabular}{lll}
\toprule
Primitive & Quantum algorithm & Impact \\
\midrule
RSA & Shor & broken (any key size) \\
Diffie--Hellman, ECDH & Shor & broken \\
ECDSA & Shor & broken \\
AES-128 & Grover & $\approx 64$-bit; use AES-256 \\
SHA-256 preimage & Grover & generic cost $\approx 2^{128}$ \\
SHA-256 collision & quantum collision algorithms & analyse separately \\
Lattice / code / hash problems & (none known) & survive \\
\bottomrule
\end{tabular}
\caption{How a large quantum computer affects common primitives. Only the number-theoretic public-key schemes are broken outright.}
\end{table}

\section{Harvest now, decrypt later}

A natural objection is that large quantum computers do not yet exist, so why migrate today? The answer is \emph{harvest now, decrypt later}. An adversary can record encrypted traffic today and decrypt it years later once a quantum computer is available. Any data that must remain confidential beyond the arrival of quantum computing is therefore already at risk.

Signatures create a different risk. A future quantum adversary may forge \emph{new} signatures, including signatures on malicious firmware or certificates, and may undermine long-term verification. Whether a historically verified signature remains useful depends on its timestamp, archival evidence, and trust model; it is not a case of ``record now, decrypt later.'' Both risks argue for migration before the threat materializes.

\section{The response: post-quantum standards}

Because Shor's algorithm targets the specific structure of factoring and discrete logarithms, the defense is to build public-key cryptography on entirely different hard problems -- ones with no known efficient quantum algorithm. After a multi-year public competition, NIST selected the first post-quantum schemes:

\begin{itemize}
  \item \textbf{FIPS 203 -- ML-KEM} (Kyber): key encapsulation, based on module lattices~\cite{fips203};
  \item \textbf{FIPS 204 -- ML-DSA} (Dilithium): signatures, based on module lattices~\cite{fips204};
  \item \textbf{FIPS 205 -- SLH-DSA} (SPHINCS\textsuperscript{+}): signatures, based only on hash functions~\cite{fips205};
  \item a future FN-DSA standard based on Falcon is tracked separately; any implementation discussion must name the exact draft or specification version it follows.
\end{itemize}

Post-quantum cryptography is commonly organized into \emph{five broad mathematical families} of hard problems, each offering a different route to quantum-resistant public-key cryptography~\cite{bernstein2009}. This is a teaching classification, not a claim that the families are cryptographically independent.

\begin{itemize}
  \item \textbf{Lattice-based} -- short vectors and noisy linear systems in high-dimensional lattices~\cite{micciancioregev2009,peikert2014}. The workhorse of the standards, balancing speed and size; the main subject of this book.
  \item \textbf{Hash-based} -- security reduced to nothing but a hash function; the most conservative choice, resting on the longest-understood assumption (Chapter~\ref{ch:slhdsa}).
  \item \textbf{Code-based} -- decoding random-looking linear codes; the oldest family, with a decades-long track record (Chapter~\ref{ch:mceliece}).
  \item \textbf{Multivariate} -- solving systems of multivariate quadratic equations over a finite field~\cite{ding2006}; very compact signatures and large keys, with a turbulent cryptanalytic history (Chapter~\ref{ch:multivariate}).
  \item \textbf{Isogeny-based} -- walking between elliptic curves along secret isogenies; the smallest keys of all, but the youngest and least settled family (Chapter~\ref{ch:isogeny}).
\end{itemize}

The NIST standards are concrete \emph{implementations} of the first two families: three are lattice schemes and one is hash-based. The code-based, multivariate, and isogeny families have no FIPS standard yet, but each offers active candidates and, crucially, a structurally different hedge should a future attack ever undermine lattices. Table~\ref{tab:pqc-families} places every standard in its family.

\begin{table}[H]
\centering
\small
\begin{tabular}{llll}
\toprule
Family & Hard problem & Standard / leading scheme & Role \\
\midrule
Lattice & Module-LWE & FIPS 203 ML-KEM & KEM \\
Lattice & Module-LWE / SIS & FIPS 204 ML-DSA & signature \\
Lattice & NTRU & FN-DSA / Falcon specifications & signature \\
Hash & preimage / collision & FIPS 205 SLH-DSA & signature \\
Code & syndrome decoding & Classic McEliece, HQC & KEM \\
Multivariate & MQ (quadratic systems) & UOV, MAYO (candidates) & signature \\
Isogeny & supersingular isogeny path & SQIsign (candidate) & signature \\
\bottomrule
\end{tabular}
\caption{Five commonly discussed mathematical families of post-quantum cryptography. FIPS 203--205 cover lattice- and hash-based schemes; other families remain active areas of research and standardization.}
\label{tab:pqc-families}
\end{table}

\section{Summary}

\begin{itemize}
  \item Today's public-key cryptography (RSA, Diffie--Hellman, ECC) rests on factoring and discrete logarithms.
  \item \textbf{Shor's algorithm} solves both in polynomial time on a quantum computer, breaking these schemes at \emph{any} key size.
  \item \textbf{Grover's algorithm} gives a quadratic generic-search speedup; key search, hash preimages, and hash collisions need separate security analyses.
  \item Confidentiality has a ``harvest now, decrypt later'' risk; signature migration concerns future forgery and long-term verification.
  \item The rest of the book surveys five broad mathematical families -- lattice, hash, code, multivariate, and isogeny -- with FIPS 203--205 as published NIST standards at this snapshot.
\end{itemize}

The remaining chapters start from the ground up: the mathematical objects -- finite fields, vectors, polynomials, and lattices -- on which these new hard problems are built.
